\documentclass[11pt,a4paper]{article}

\usepackage[a4paper, margin=2cm]{geometry}
\usepackage{amsmath, amssymb}
\usepackage{booktabs}
\usepackage{xcolor}
\usepackage{longtable}
\usepackage{array}
\usepackage{parskip}
\usepackage{titlesec}
\usepackage{mdframed}

\definecolor{incgreen}{RGB}{0,140,0}
\definecolor{exred}{RGB}{180,0,0}
\definecolor{borderorange}{RGB}{200,120,0}

\newcommand{\inc}{\textbf{\textcolor{incgreen}{INCLUDE}}}
\newcommand{\exc}{\textbf{\textcolor{exred}{EXCLUDE}}}
\newcommand{\bord}{\textbf{\textcolor{borderorange}{BORDERLINE}}}

\titleformat{\section}{\large\bfseries}{}{0em}{}[\titlerule]
\titleformat{\subsection}{\normalsize\bfseries}{}{0em}{}

\pagestyle{plain}

\begin{document}

\begin{center}
  {\LARGE \textbf{Mathematics for Physics 2 --- Week 1 Review}} \\[6pt]
  {\large Sources: \textit{Memphis\_2\_week\_1.pdf}, \textit{ClassEx1 (2026).pdf}} \\[4pt]
  {\normalsize Pending approval before writing to cheat sheet}
\end{center}

\vspace{1em}

\begin{mdframed}[backgroundcolor=gray!10]
\textbf{Legend:} \inc\ --- will be added to cheat sheet \quad
\exc\ --- excluded \quad \bord\ --- needs your decision
\end{mdframed}

\vspace{1em}

% ------------------------------------------------------------------
\section{Notation \& Definitions}

\begin{longtable}{p{0.5cm} p{7.5cm} p{2.2cm} p{3.5cm}}
\toprule
\# & Formula / Definition & Decision & Reason \\
\midrule
\endhead

N1 & \textbf{Hat notation:} $\hat{v}$ denotes a \emph{unit vector} in the direction of $\vec{v}$:
\[
  \hat{v} = \frac{\vec{v}}{|\vec{v}|}, \qquad |\hat{v}| = 1
\]
& \inc & Foundational notation \\[14pt]

N2 & \textbf{Vector components:} A vector $\vec{A}$ in 3D Cartesian coordinates is written
\[
  \vec{A} = A_x\hat{x} + A_y\hat{y} + A_z\hat{z}
\]
$A_i$ denotes the $i$-th component of $\vec{A}$ (e.g.\ $A_1 = A_x$). \newline
The magnitude is $|\vec{A}| = A = \sqrt{A_x^2 + A_y^2 + A_z^2}$.
& \inc & Required to understand index notation \\[14pt]

\bottomrule
\end{longtable}

% ------------------------------------------------------------------
\section{Index Notation \& Einstein Summation}

\begin{longtable}{p{0.5cm} p{7.5cm} p{2.2cm} p{3.5cm}}
\toprule
\# & Formula & Decision & Reason \\
\midrule
\endhead
1 & $\vec{A} = A_i\hat{e}_i$ \quad (sum over $i=1,2,3$ implied) & \inc & Core notation \\[10pt]
2 & \textbf{Einstein convention:} any index repeated \emph{exactly twice} in a term implies summation over 1, 2, 3. Example:
\[a_i b_i \;\equiv\; \sum_{i=1}^{3} a_i b_i\]
& \inc & Fundamental rule \\[14pt]
3 & \textbf{Index conservation:} free indices must be the same on both sides of an equation. & \inc & Key consistency rule \\[6pt]
\bottomrule
\end{longtable}

% ------------------------------------------------------------------
\section{Kronecker Delta}

\begin{longtable}{p{0.5cm} p{7.5cm} p{2.2cm} p{3.5cm}}
\toprule
\# & Formula & Decision & Reason \\
\midrule
\endhead
4 & $\displaystyle\delta_{ij} = \begin{cases} 1 & i = j \\ 0 & i \neq j \end{cases}$ & \inc & Definition \\[18pt]
5 & $\delta_{ij} = \delta_{ji}$ & \inc & Symmetry \\[10pt]
6 & $A_j\,\delta_{ij} = A_i$ & \inc & Index contraction --- used constantly \\[10pt]
7 & $\delta_{ik}\,\delta_{jk} = \delta_{ij}$ & \inc & Product identity \\[10pt]
\textit{(was 8)} & $\delta_{ii} = 3$ & \exc & \textit{Dropped per review} \\[6pt]
\bottomrule
\end{longtable}

% ------------------------------------------------------------------
\section{Levi-Civita Symbol}

\begin{longtable}{p{0.5cm} p{7.5cm} p{2.2cm} p{3.5cm}}
\toprule
\# & Formula & Decision & Reason \\
\midrule
\endhead
10 & $\displaystyle\varepsilon_{ijk} = \begin{cases} +1 & (i,j,k)\ \text{cyclic: } (1,2,3),(2,3,1),(3,1,2) \\ -1 & \text{anti-cyclic: } (2,1,3),(1,3,2),(3,2,1) \\ 0 & \text{any two indices equal} \end{cases}$ & \inc & Definition \\[26pt]
11 & $\varepsilon_{ijk} = -\varepsilon_{jik}$ \quad (swapping any two adjacent indices flips sign) & \inc & Key property \\[10pt]
\textit{(was 12)} & $\varepsilon_{ijk} = \varepsilon_{kij} = \varepsilon_{jki}$ & \exc & \textit{Dropped per review} \\[10pt]
13 & $\displaystyle\sum_k \varepsilon_{ijk}\,\varepsilon_{mnk} = \delta_{im}\delta_{jn} - \delta_{in}\delta_{jm}$ & \inc & Simplifies double cross products \\[6pt]
\bottomrule
\end{longtable}

% ------------------------------------------------------------------
\section{Dot Product (Scalar Product)}

\begin{longtable}{p{0.5cm} p{7.5cm} p{2.2cm} p{3.5cm}}
\toprule
\# & Formula & Decision & Reason \\
\midrule
\endhead
14 & $\vec{A}\cdot\vec{B} = AB\cos\theta$ \quad ($\theta$ = angle between $\vec{A}$ and $\vec{B}$) & \inc & Geometric definition \\[10pt]
15 & $\vec{A}\cdot\vec{B} = A_i B_i = \displaystyle\sum_{i=1}^{3} A_i B_i$ & \inc & Index form \\[14pt]
16 & $\hat{e}_i\cdot\hat{e}_j = \delta_{ij}$ & \inc & Orthonormality \\[10pt]
17 & $|\vec{A}| = \sqrt{\vec{A}\cdot\vec{A}} = \sqrt{A_i A_i} = \sqrt{\displaystyle\sum_{i=1}^{3} A_i^2}$\newline
\textit{Note: $A_i A_i$ uses Einstein convention --- $i$ appears twice, so it sums: $A_1^2 + A_2^2 + A_3^2$.}
& \inc & Clarified \\[18pt]
18 & $\vec{A}\cdot\vec{B}=0 \Rightarrow \vec{A}\perp\vec{B}$ & \exc & Trivially obvious from \#14 \\[10pt]
19 & $\vec{A}\cdot\vec{B}=AB$ when parallel & \exc & Special case of \#14 \\[10pt]
20 & $\vec{A}\cdot\vec{B}=\vec{B}\cdot\vec{A}$ & \exc & Obvious from $A_iB_i = B_iA_i$ \\[6pt]
\bottomrule
\end{longtable}

% ------------------------------------------------------------------
\section{Cross Product (Vector Product)}

\begin{longtable}{p{0.5cm} p{7.5cm} p{2.2cm} p{3.5cm}}
\toprule
\# & Formula & Decision & Reason \\
\midrule
\endhead
21 & $|\vec{a}\times\vec{b}| = ab\sin\theta$ \quad ($\theta$ = angle between $\vec{a}$ and $\vec{b}$) & \inc & Geometric definition \\[10pt]
22 & If $\vec{C} = \vec{A}\times\vec{B}$, then the $i$-th component of $\vec{C}$ is:
\[ C_i = \varepsilon_{ijk}\,A_j B_k \]
($j$ and $k$ are summed; $i$ is the free index selecting the component)
& \inc & Clarified \\[18pt]
23 & $\vec{a}\times\vec{b} = \begin{vmatrix} \hat{x} & \hat{y} & \hat{z} \\ a_x & a_y & a_z \\ b_x & b_y & b_z \end{vmatrix}$ & \inc & Practical computation formula \\[24pt]
24 & $(a_yb_z - a_zb_y)\hat{x} + (a_zb_x - a_xb_z)\hat{y} + (a_xb_y - a_yb_x)\hat{z}$ & \exc & Derivable from \#23 \\[10pt]
25 & $\vec{a}\times\vec{b} = -\vec{b}\times\vec{a}$ & \inc & Anti-commutativity \\[10pt]
26 & $\hat{x}\times\hat{y}=\hat{z},\quad \hat{y}\times\hat{z}=\hat{x},\quad \hat{z}\times\hat{x}=\hat{y}$ \newline
   $\hat{y}\times\hat{x}=-\hat{z},\quad \hat{z}\times\hat{y}=-\hat{x},\quad \hat{x}\times\hat{z}=-\hat{y}$ \newline
   $\hat{x}\times\hat{x}=\hat{y}\times\hat{y}=\hat{z}\times\hat{z}=\vec{0}$ & \inc & Quick calculation \\[10pt]
27 & $\vec{a}\cdot(\vec{a}\times\vec{b})=0$ & \exc & Consequence of $\varepsilon$ \\[6pt]
\bottomrule
\end{longtable}

% ------------------------------------------------------------------
\section{Key Vector Identity}

\begin{longtable}{p{0.5cm} p{7.5cm} p{2.2cm} p{3.5cm}}
\toprule
\# & Formula & Decision & Reason \\
\midrule
\endhead
28 & \textbf{BAC-CAB rule:}
\[ \vec{A}\times(\vec{B}\times\vec{C}) = \vec{B}\,(\vec{A}\cdot\vec{C}) - \vec{C}\,(\vec{A}\cdot\vec{B}) \]
Here $(\vec{A}\cdot\vec{C})$ and $(\vec{A}\cdot\vec{B})$ are \emph{scalars} (plain numbers); multiplying a scalar by a vector $\vec{B}$ or $\vec{C}$ scales its magnitude. & \inc & Clarified \\[14pt]
\bottomrule
\end{longtable}

% ------------------------------------------------------------------
\section{Coordinate Systems}

\subsection{Polar Coordinates $(r, \theta)$}

\begin{longtable}{p{0.5cm} p{7.5cm} p{2.2cm} p{3.5cm}}
\toprule
\# & Formula & Decision & Reason \\
\midrule
\endhead
29 & $x = r\cos\theta,\quad y = r\sin\theta$ \newline
   $r = \sqrt{x^2+y^2},\quad \theta = \arctan\!\left(\dfrac{y}{x}\right)$ & \inc & Conversion formulas \\[14pt]
30 & $\hat{r} = \cos\theta\,\hat{x} + \sin\theta\,\hat{y}$ \newline
   $\hat{\theta} = -\sin\theta\,\hat{x} + \cos\theta\,\hat{y}$ & \inc & Unit vectors in polar \\[10pt]
31 & $\hat{x} = \cos\theta\,\hat{r} - \sin\theta\,\hat{\theta}$ \newline
   $\hat{y} = \sin\theta\,\hat{r} + \cos\theta\,\hat{\theta}$ & \inc & Inverse conversion \\[6pt]
\bottomrule
\end{longtable}

\subsection{Spherical Coordinates $(r, \theta, \varphi)$}

\begin{longtable}{p{0.5cm} p{7.5cm} p{2.2cm} p{3.5cm}}
\toprule
\# & Formula & Decision & Reason \\
\midrule
\endhead
32 & $x = r\sin\theta\cos\varphi,\quad y = r\sin\theta\sin\varphi,\quad z = r\cos\theta$ \newline
   $r = \sqrt{x^2+y^2+z^2},\quad \varphi = \arctan\!\left(\dfrac{y}{x}\right),\quad \theta = \arccos\!\left(\dfrac{z}{r}\right)$ & \inc & Conversion formulas \\[14pt]
33 & $\hat{r} = \sin\theta\cos\varphi\,\hat{x} + \sin\theta\sin\varphi\,\hat{y} + \cos\theta\,\hat{z}$ \newline
   $\hat{\theta} = \cos\theta\cos\varphi\,\hat{x} + \cos\theta\sin\varphi\,\hat{y} - \sin\theta\,\hat{z}$ \newline
   $\hat{\varphi} = -\sin\varphi\,\hat{x} + \cos\varphi\,\hat{y}$ & \inc & Unit vectors in spherical \\[14pt]
34 & $\hat{x} = \sin\theta\cos\varphi\,\hat{r} + \cos\theta\cos\varphi\,\hat{\theta} - \sin\varphi\,\hat{\varphi}$ \newline
   $\hat{y} = \sin\theta\sin\varphi\,\hat{r} + \cos\theta\sin\varphi\,\hat{\theta} + \cos\varphi\,\hat{\varphi}$ \newline
   $\hat{z} = \cos\theta\,\hat{r} - \sin\theta\,\hat{\theta}$ & \inc & Included per review \\[6pt]
\bottomrule
\end{longtable}

\end{document}
